% Default CV adapted from the CERN STAG 2026 PDF layout, 24 October 2026.
% NATO bullets selected from Shared/Professional Experience/NATO CMRE evidence.
% Graduation months and project context updated at the candidate's request.
% Google coursework review, 24 October 2026: candidate confirms taking all
% courses in the supplied SOIA and UFRGS curricula. English course labels
% translate those PDFs; selected for ML/research and software engineering.
% Build from this directory: latexmk -pdf -outdir=/tmp/default-cv-build Rafael_Ode_Brino_CV.tex

\DocumentMetadata{lang=en-US,pdfversion=1.7,tagging=on}
\documentclass[a4paper,11pt]{article}
\usepackage[T1]{fontenc}
\usepackage[utf8]{inputenc}
\usepackage{lmodern}
\usepackage[a4paper,left=13mm,right=13mm,top=12mm,bottom=13mm]{geometry}
\usepackage[hidelinks]{hyperref}
\usepackage[american]{babel}
\usepackage{microtype}
\input{glyphtounicode}
\pdfgentounicode=1

% Snapshot of contact and degree values for a self-contained, editable CV.
% Shared rendering values used by the CV and CL templates.
% Factual reference: Profile.md. Keep these values consistent with that note.
\newcommand{\CandidateName}{Rafael Ode Brino}
\newcommand{\CandidateLocation}{Brest, France}
\newcommand{\CandidatePhone}{+33 7 45 75 03 71}
\newcommand{\CandidatePhoneURL}{tel:+33745750371}
\newcommand{\CandidateEmail}{rafael.ode_brino@ensta.fr}
\newcommand{\CandidateEmailText}{rafael.ode\_brino@ensta.fr}
\newcommand{\CandidateLinkedInURL}{https://linkedin.com/in/odebrino}
\newcommand{\CandidateLinkedInText}{linkedin.com/in/odebrino}
\newcommand{\CandidateGitHubURL}{https://github.com/odebrino}
\newcommand{\CandidateGitHubText}{github.com/odebrino}
\newcommand{\CandidateCitizenship}{German \& Brazilian}
\newcommand{\CandidateAvailability}{a six-month ENSTA final-year research internship starting in March 2027}

% Shared rendering values. Factual reference: Education.md.
% Programme PDFs are curricula, not personal transcripts.
\newcommand{\CandidateSchool}{ENSTA}
\newcommand{\CandidateSchoolGroup}{Institut Polytechnique de Paris}
\newcommand{\CandidateSchoolLocation}{Brest, France}
\newcommand{\CandidateSchoolDescription}{French graduate engineering school}
\newcommand{\CandidateDegree}{MSc in Engineering}
\newcommand{\CandidateSpecialization}{Observation Systems \& Artificial Intelligence}
\newcommand{\CandidateSpecializationProse}{Observation Systems and Artificial Intelligence}
\newcommand{\CandidateSchoolDates}{2025--2027}
\newcommand{\CandidateSchoolGraduation}{September 2027}
\newcommand{\CandidateBachelorsSchool}{Federal University of Rio Grande do Sul (UFRGS)}
\newcommand{\CandidateBachelorsLocation}{Porto Alegre, Brazil}
\newcommand{\CandidateBachelorsDegree}{BSc in Electrical Engineering}
\newcommand{\CandidateBachelorsDates}{2023--2028}
\newcommand{\CandidateBachelorsGraduation}{August 2028}


% Match the supplied CV: A4, serif type, one column, ruled headings.
\pagestyle{empty}
\urlstyle{same}
\raggedbottom
\raggedright
\setlength{\tabcolsep}{0in}
% Use a consistent baseline, then separate titles, bullets and sections.
\makeatletter
\renewcommand\small{\@setfontsize\small{10.2}{13.6}}
\makeatother

\hypersetup{
  pdftitle={\CandidateName\space - Curriculum Vitae},
  pdfauthor={\CandidateName},
  pdfsubject={Machine Learning and Scientific Computing}
}

% Native sectioning supports structured PDF headings and avoids titlesec's
% incompatibility with automatic tagging. Keep the same restrained ruled style.
\makeatletter
\renewcommand\section{\@startsection{section}{1}{0pt}{9pt}{3pt}%
  {\normalfont\large\scshape\raggedright}}
\makeatother
\newcommand{\cvSection}[1]{%
  \section*{#1}\nobreak\vspace{-2pt}\hrule\nobreak\vspace{6pt}}

% Layout helpers. Edit content below rather than spacing in individual entries.
\newcommand{\cvItem}[1]{\item #1\par}
\newcommand{\cvEntry}[4]{%
  \item
  \textbf{#1}\hfill #2\par\nobreak
  {\small\textit{#3}\hfill\textit{#4}\par}
}
\newcommand{\cvEducationEntry}[5]{%
  \item
  \textbf{#1}\hfill #2\par\nobreak
  {\small\textit{#3}\hfill\textit{#4}\par\nobreak
  #5\par}
}
\newcommand{\cvProject}[2]{%
  \item
  {\small #1\hfill #2\par}
}
\renewcommand\labelitemii{\textbullet}
\newenvironment{cvEntries}
  {\begin{itemize}[left-margin=0.15in,item-label={},begin-vspace=0pt,
    begin-extra-vspace=0pt,item-vspace=6pt,para-vspace=0pt]}{\end{itemize}}
\newenvironment{cvBullets}
  {\small\begin{itemize}[left-margin=0.18in,begin-vspace=3pt,
    begin-extra-vspace=0pt,item-vspace=3pt,para-vspace=0pt]}{\end{itemize}}

% Each manuscript is a separate group. Its bullets sit to the right of its
% heading; addvspace supplies a visible break before the next manuscript.
% Reset cvarticle to zero when adding a different employer's manuscripts.
\newcounter{cvarticle}
\newenvironment{cvArticle}[2][in preparation for submission]{%
  \par\addvspace{6pt}\stepcounter{cvarticle}%
  {\small\hspace*{0.12in}\textbf{Article \arabic{cvarticle}: #2}
    \textit{(#1)}\par}%
  \small\begin{itemize}[left-margin=0.34in,item-label=\textbullet,
    begin-vspace=3pt,item-vspace=3pt,para-vspace=0pt]
}{\end{itemize}}

% Project-sourced skills: one spaced list item per category.
\newenvironment{cvSkills}
  {\small\begin{itemize}[left-margin=0.15in,item-label={},begin-vspace=2pt,
    begin-extra-vspace=0pt,item-vspace=2pt,para-vspace=0pt]}{\end{itemize}}
\newcommand{\cvSkill}[2]{\item \textbf{#1:} #2}

\begin{document}

% CONTACT ---------------------------------------------------------------
\begin{center}
  \textbf{\huge \CandidateName} \\ \vspace{2pt}
  \small \CandidateLocation{} \textbar{} \CandidateCitizenship{} citizenship \\ \vspace{2pt}
  \small
  \href{\CandidatePhoneURL}{\CandidatePhone} \textbar{}
  \href{mailto:\CandidateEmail}{\CandidateEmailText} \textbar{}
  \href{\CandidateLinkedInURL}{\CandidateLinkedInText} \textbar{}
  \href{\CandidateGitHubURL}{\CandidateGitHubText}
\vspace{2pt}\\
  \small Six-month ENSTA final-year research internship starting in March 2027
\end{center}

\cvSection{Education}
\begin{cvEntries}
  \cvEducationEntry
    {\CandidateSchool{} --- \CandidateSchoolGroup}{\CandidateSchoolLocation}
    {\CandidateDegree{} --- \CandidateSpecialization; expected \CandidateSchoolGraduation}{\CandidateSchoolDates}
    {Selective French engineering school (\textit{grande \'ecole}); concurrent studies with UFRGS.}
  \begin{cvBullets}
    \cvItem{Relevant coursework: Machine Learning, Deep Learning, Numerical Optimization \& Markov Models, Digital Image Processing, Software Design, and Decision Theory \& Estimation.}
  \end{cvBullets}

  \cvEntry
    {\CandidateBachelorsSchool}{\CandidateBachelorsLocation}
    {\CandidateBachelorsDegree; expected \CandidateBachelorsGraduation}{\CandidateBachelorsDates}
  \begin{cvBullets}
    \cvItem{Relevant coursework: Data Structures, Algorithms \& Programming, Object-Oriented Programming, Probability \& Statistics, Linear Algebra, and Numerical Calculus.}
  \end{cvBullets}
\end{cvEntries}

\cvSection{Research Experience}
\begin{cvEntries}
\cvEntry{NATO Science \& Technology Organization (STO)}{La Spezia, Italy}
{Research Intern --- Centre for Maritime Research and Experimentation (CMRE)}{May--September 2026}
\setcounter{cvarticle}{0}
% CMRE evidence: G-H support the Bayesian implementation and marginalization;
% I-J support group-averaged AUC, bootstrap ranking gains, and the absence
% of stable gains from richer physical-quality/nonlinear specifications.
% A/E support the delivered pipeline and 64 paired experimental conditions;
% F supports the 48.45% response-prediction RMSE reduction under nested selection.
% Candidate confirms PyTorch/TensorFlow and Slurm/HPC use at NATO,
% 24 October 2026, in response to the training/inference/compute-workflow query.
\begin{cvArticle}[first author; in preparation]{Hierarchical Bayesian Modeling of Detection Success}
  \cvItem{Implemented \textbf{hierarchical Bayesian probit models} in Python for object-detection success, combining partial pooling and MCMC with analytical marginalization for predictions in unseen data-collection groups.}
  \cvItem{Achieved \textbf{group-averaged ROC AUC 0.734} on 21 held-out groups; nested validation and paired group bootstraps supported ranking gains over a prevalence baseline, while richer features added no stable benefit.}
\end{cvArticle}
\begin{cvArticle}[first author; in preparation]{Modeling Detector Performance Under Image Degradation}
  \cvItem{Built a Python/PyTorch pipeline for YOLO training and evaluation on Slurm-managed HPC; ran paired tests across \textbf{64 conditions} of resolution, blur, and noise, delivered reusable code, and presented findings at CMRE.}
  \cvItem{Calibrated a nonlinear model of detector performance, reducing prediction \textbf{RMSE by 48\% (\mbox{0.237 to 0.122})} versus reusing the source coefficients unchanged; evaluated with fully nested selection on held-out noise levels.}
\end{cvArticle}
\end{cvEntries}
\cvSection{Selected Projects}
\begin{cvEntries}
\cvProject{\textbf{FX Risk Forecasting} \textbar{} \emph{Python, C++, SQL, DuckDB}}{2026--Present}
\begin{cvBullets}
  \cvItem{Built a Python/SQL (DuckDB) pipeline for ECB exchange rates and implemented HMM, Student-t GARCH, and independent-mixture models to test whether volatility regimes improve forecasts of currency-return distributions.}
  \cvItem{Implemented a \textbf{C++ HMM filter} with Python bindings and numerical parity checks; built rolling out-of-sample evaluation with training-cutoff checks and paired block-bootstrap comparisons.}
\end{cvBullets}
\cvProject{\textbf{Augmented Reality for Surgical Training} \textbar{} \emph{Python, C++, C\#, Unity}}{2025--2026}
\begin{cvBullets}
  \cvItem{Built a computer-vision pipeline for HRIACT -- French Armed Forces hospital, mapping \textbf{physical scalpel movements} from Python/C++ to a virtual instrument in a Unity/C\# surgical-training simulation.}
\end{cvBullets}
\cvProject{\textbf{3D Mapping with Detection and SLAM} \textbar{} \emph{Python, YOLO}}{2024--2025}
\begin{cvBullets}
  \cvItem{Co-developed an EKF-SLAM pipeline fusing stereo visual odometry, YOLO detections, and GPS for 3D mapping and trajectory estimation; coauthored the \textbf{published CBA 2026 preprint} (\href{https://ojs.sba.org.br/index.php/cba/article/download/6678/6025}{\underline{paper; English abstract}}).}
\end{cvBullets}
\end{cvEntries}
\cvSection{Technical Skills}
% Technical stack selected from Shared evidence and the candidate's explicit
% confirmation of Git/GitHub/GitLab, Bash/shell, pytest, JupyterLab,
% PyTorch/TensorFlow, MATLAB/Simulink, LaTeX, SSH/Slurm/HPC and bindings,
% 24 October 2026. Linux and Neovim were confirmed in the preceding request.
% Select pybind11 from the confirmed binding tools; no project-specific
% attribution, measured speedup or proficiency level is inferred.
% Compact methods row restored from CMRE/HMM, sensor-fusion methods and
% candidate-confirmed numerical-optimization coursework. No proficiency scale.
% SLAM publication link verified on the official SBA CBA page, 24 October 2026:
% https://ojs.sba.org.br/index.php/cba/article/view/6678 (Rafael O. Brino listed).
\begin{cvSkills}
  \cvSkill{Programming}{Python, C++, SQL, Bash, MATLAB, C\#}
  \cvSkill{ML \& Data}{PyTorch, TensorFlow, scikit-learn, NumPy, SciPy, pandas, OpenCV, DuckDB}
  \cvSkill{Modeling \& Inference}{Bayesian inference, MCMC, HMMs, EM, nonlinear optimization, Kalman filtering}
  \cvSkill{Development \& HPC}{Git (GitHub/GitLab), Linux, pytest, pybind11, SSH, Slurm, Neovim}
  \cvSkill{Research \& Simulation}{JupyterLab, Unity, Simulink, Matplotlib, LaTeX}
\end{cvSkills}

\cvSection{Languages}
\begin{itemize}[left-margin=0.15in,item-label={}]
  \small\item{Portuguese: Native \textbar{} English: Fluent \textbar{} French: Advanced \textbar{} German: Elementary}
\end{itemize}
\end{document}
